% ============================================================
%   Chapter 5  \textemdash\  Implementation
% ============================================================

\chapter{Implementation}\label{ch:implementation}

This chapter describes how the design of Chapter~\ref{ch:design} is materialised in Python code.
Section~\ref{sec:impl-stack} introduces the technology stack and
Section~\ref{sec:impl-layout} the project layout. The subsequent sections cover each component:
the CARLA wrapper (\S\ref{sec:impl-env}), the PPO agent (\S\ref{sec:impl-agent}), the two shields
(\S\ref{sec:impl-shields}), the reward shaper (\S\ref{sec:impl-reward}), the curriculum manager
(\S\ref{sec:impl-curriculum}), the entry-point pipelines (\S\ref{sec:impl-pipelines}) and
reproducibility concerns (\S\ref{sec:impl-repro}).

\section{Technology stack}\label{sec:impl-stack}

The framework is written in pure Python (\(\ge 3.10\)) and relies on the libraries summarised in
Table~\ref{tab:stack}. Licences are reported because the framework is meant to be distributed as
open academic software.

\begin{table}[h]
  \centering
  \caption{Main third-party dependencies.}
  \label{tab:stack}
  \small
  \begin{tabular}{@{}lll@{}}
    \toprule
    \textbf{Library} & \textbf{Role}                              & \textbf{Licence} \\
    \midrule
    \texttt{carla}         & Python binding for CARLA 0.9.16            & MIT \\
    \texttt{gymnasium}     & Standard RL environment API                & MIT \\
    \texttt{torch}         & Deep learning backend                      & BSD-3 \\
    \texttt{numpy}         & Numerical arrays                           & BSD-3 \\
    \texttt{matplotlib}    & Evaluation dashboard                       & PSF-based \\
    \texttt{streamlit}     & Live training dashboard                    & Apache-2.0 \\
    \texttt{plotly}        & Interactive charts in the dashboard        & MIT \\
    \texttt{sqlite3} (stdlib) & Metrics database                        & Public domain \\
    \texttt{ruff}          & Linter and formatter                       & MIT \\
    \texttt{pytest}        & Unit testing of the reward shaper          & MIT \\
    \bottomrule
  \end{tabular}
\end{table}

\section{Project layout}\label{sec:impl-layout}

The repository root is organised as follows (directories without relevance to the thesis are
omitted):

\begin{itemize}
  \item \texttt{main\_train.py} \textemdash\ training entry point.
  \item \texttt{main\_eval.py} \textemdash\ evaluation entry point.
  \item \texttt{dataVisualizer.py} \textemdash\ Streamlit dashboard that reads the SQLite
        metrics database.
  \item \texttt{src/CARLA/Env/carla\_env.py} \textemdash\ base Gymnasium environment.
  \item \texttt{src/PPO/} \textemdash\ \texttt{ActorCritic.py}, \texttt{ppo\_agent.py},
        \texttt{RunningMeanStd.py}.
  \item \texttt{src/safety\_shield.py} \textemdash\ basic safety shield.
  \item \texttt{src/Adaptative\_Shield/} \textemdash\ \texttt{adaptive\_horizon\_shield.py},
        \texttt{BicycleModel.py}.
  \item \texttt{src/reward\_shaper.py} \textemdash\ dense reward shaper.
  \item \texttt{src/curriculumManager.py} \textemdash\ performance-driven curriculum.
  \item \texttt{src/Metrics/} \textemdash\ live metrics logger and offline evaluation reporter.
  \item \texttt{tests/test\_reward\_balance.py} \textemdash\ unit tests on the reward shaper.
  \item \texttt{data/models/} \textemdash\ serialised checkpoints.
  \item \texttt{runs/} \textemdash\ per-run SQLite databases and TensorBoard logs.
\end{itemize}

Every module exposes either a class that inherits from \texttt{gymnasium.Wrapper} (the shields,
the reward shaper) or a small self-contained helper (the curriculum, the running-mean estimator).
There is no inter-module state other than what flows through the Gymnasium interface, which makes
individual components trivially testable.

\section{Environment wrapper}\label{sec:impl-env}

The class \texttt{CarlaEnv} inherits from \texttt{gymnasium.Env} and encapsulates the connection
to the CARLA server. At construction time it performs the following setup:

\begin{enumerate}
  \item Opens a TCP connection on \texttt{(host, port)} and loads the requested map.
  \item Enables synchronous mode with \texttt{fixed\_delta\_seconds = 0.05}.
  \item Instantiates the Traffic Manager on \texttt{tm\_port} and spawns the requested number
        of NPC vehicles.
  \item Sets the requested weather preset.
  \item Registers the collision, lane-invasion and semantic-LIDAR sensors.
  \item Defines \texttt{observation\_space} as \(\text{Box}([0,1]^{739})\) and
        \texttt{action\_space} as \(\text{Box}([-1,1]^2)\).
\end{enumerate}

On every \texttt{step(action)} call, the class converts the normalised action into a
\texttt{carla.VehicleControl}, ticks the simulator exactly once, reads the sensors, computes the
geometric features from the Waypoint API, builds the 739-dimensional observation, and derives the
base reward together with the \texttt{done} and \texttt{truncated} flags. The \texttt{info}
dictionary contains every scalar used by downstream wrappers.

\section{PPO agent}\label{sec:impl-agent}

The PPO agent is split into two files: \texttt{src/PPO/ActorCritic.py} contains the network
described in Figure~\ref{fig:actor-critic-arch}, and \texttt{src/PPO/ppo\_agent.py} contains
the training loop. The three public methods used by \texttt{main\_train.py} are
\texttt{select\_action}, \texttt{evaluate\_action} and \texttt{update}.
Figure~\ref{fig:select-action} shows the pseudocode of \texttt{select\_action} and
\texttt{evaluate\_action}; Figure~\ref{fig:ppo-update} shows the pseudocode of \texttt{update}.

\begin{figure}[h]
  \centering
  \includegraphics[width=0.95\linewidth]{imgs/code/select_action}
  \caption{Action selection and evaluation methods. \texttt{select\_action} normalises the
           vector block of the observation, samples from the policy (or returns the mean for
           deterministic evaluation), and returns the action, log probability and value estimate.
           \texttt{evaluate\_action} is called with \texttt{info["executed\_action"]} to
           recompute the log probability of the action that actually reached the environment.}
  \label{fig:select-action}
\end{figure}

\begin{figure}[h]
  \centering
  \includegraphics[width=0.95\linewidth]{imgs/code/ppo_update}
  \caption{PPO update loop with GAE backward pass, return normalisation by running standard
           deviation, and up to \(K=10\) epochs of clipped surrogate loss with optional
           value clipping and KL early stop.}
  \label{fig:ppo-update}
\end{figure}

\subsection{Credit assignment under a shield}\label{sec:impl-shield-credit}

The interplay between the shield and the PPO update is the single most error-prone point of the
framework. Let \(a_t\) be the action sampled from the policy and
\(a^\star_t = \Sigma(s_t, a_t)\) the action produced by the shield. The shield is a
\texttt{gymnasium.Wrapper}, so the environment sees \(a^\star_t\); the agent only receives it
through the \texttt{executed\_action} key of the \texttt{info} dictionary.

The training loop in \texttt{main\_train.py} stores \(a^\star_t\) (not \(a_t\)) in the rollout
buffer and recomputes the log probability of \(a^\star_t\) under the current policy via
\texttt{evaluate\_action}. Without this step, the gradient estimator would update the policy
towards the action that \emph{was not executed}, effectively decoupling the policy from the
environment and making training degenerate. With this step, the policy learns to produce actions
that are close to what the shield considered safe, which is precisely the desired behaviour: the
shield teaches the agent what \emph{not} to do.

The approximate KL used by the early-stop mechanism (Figure~\ref{fig:ppo-update}) is the
unbiased estimator of~\cite{schulman2020approxkl}. The KL target and the value-clip parameter
\(\xi\) are both exposed as command-line flags.

\section{Shield implementations}\label{sec:impl-shields}

Both shields inherit from \texttt{gymnasium.Wrapper} and implement the same public interface
(\texttt{reset}, \texttt{step}, \texttt{get\_statistics}, \texttt{reset\_statistics}). The
interchangeability is what makes \texttt{--shield\_type} a drop-in parameter.

\subsection{Basic shield}\label{sec:impl-basic}

The implementation of \texttt{CarlaSafetyShield} follows the correction law of
Section~\ref{sec:design-basic-shield} with steering gain \(k_s = 2.5\) and brake gain
\(k_b = 8.0\). The class keeps per-episode statistics of the sub-reason that triggered each
intervention (\texttt{front}, \texttt{side\_right}, \texttt{side\_left}, \texttt{lane\_right},
\texttt{lane\_left}, \texttt{heading}), surfaced by \texttt{get\_statistics()} and logged once
per episode.

\subsection{Adaptive-horizon shield}\label{sec:impl-adaptive}

The adaptive shield is more involved. Figure~\ref{fig:adaptive-step} shows the structure of its
\texttt{step} method, which classifies the current risk level, queries the bicycle model for
the appropriate horizon, checks the proposed action, searches for a safe alternative if needed,
and decorates the \texttt{info} dictionary before returning.

\begin{figure}[h]
  \centering
  \includegraphics[width=0.95\linewidth]{imgs/code/adaptive_shield_step}
  \caption{Adaptive-horizon shield \texttt{step} method. The semantic analysis extracts LIDAR
           and Waypoint features; the risk classifier selects the horizon \(H\); the trajectory
           check decides whether to pass the proposed action through or invoke the candidate
           search; the \texttt{info} dictionary is updated with shield diagnostics.}
  \label{fig:adaptive-step}
\end{figure}

Figure~\ref{fig:adaptive-check} shows the trajectory-check routine: it validates LIDAR
thresholds, checks for a pedestrian emergency, and then runs the bicycle model over \(H\) steps,
projecting each predicted position onto the road network via CARLA's Waypoint API.

\begin{figure}[h]
  \centering
  \includegraphics[width=0.95\linewidth]{imgs/code/adaptive_trajectory_check}
  \caption{Trajectory safety check of the adaptive shield. The threshold multiplier \(\kappa\)
           scales the LIDAR and lateral thresholds according to the current risk level. The
           bicycle model predicts \(H\) positions; each is projected onto the nearest driving
           lane and its lateral offset is checked. The lane-change relaxation (\(\tau_l = 1.2\))
           prevents the shield from blocking legitimate manoeuvres.}
  \label{fig:adaptive-check}
\end{figure}

The candidate-search routine follows the priority-ordered policy of
Section~\ref{sec:design-adaptive-shield}. Each candidate is a hard-coded two-dimensional vector;
the first that passes \texttt{\_check\_trajectory\_safety} is returned, and if none passes, the
fallback action is returned.

\subsection{Bicycle model}\label{sec:impl-bicycle}

\texttt{BicycleModel.predict\_trajectory} is a loop of \(H\) discrete steps that implements
the forward-Euler integration with closed-form arc update when \(|\delta| \ge 10^{-4}\) rad.
The class is a plain Python object with no PyTorch or CARLA dependency and is therefore
trivially unit-testable.

\section{Reward shaper}\label{sec:impl-reward}

\texttt{CarlaRewardShaper} reads approximately 20 scalars from the \texttt{info} dictionary and
produces the components of Equation~\eqref{eq:shaped-reward} from Chapter~\ref{ch:design}. It
also writes every sub-component back into \texttt{info} under clearly named keys
(\texttt{speed\_bonus}, \texttt{lane\_center\_bonus}, \texttt{heading\_bonus}, etc.) so the
metrics logger can record the per-step reward decomposition. Three minor state variables are kept
across steps:

\begin{itemize}
  \item \texttt{\_last\_steering}: needed by the smoothness penalty.
  \item \texttt{\_last\_milestone}: needed by the progress bonus to avoid double-crediting the
        same \(25\,\mathrm{m}\) segment.
  \item \texttt{\_shield\_grace\_steps}: count-down of the grace period that temporarily
        disables the idle penalty after a shield activation.
\end{itemize}

The unit test \texttt{tests/test\_reward\_balance.py} feeds a crafted \texttt{info} dictionary
to the shaper and asserts that each component matches its analytical expression. The test runs
in \(< 1\,\mathrm{s}\) without a CARLA server and is used as a pre-commit safeguard against
regressions.

\section{Curriculum manager}\label{sec:impl-curriculum}

\texttt{CurriculumManager} is a small self-contained class. Its only state is the tuple
\((i,\, E_{\mathrm{stage}},\, C_{\mathrm{rb}})\) where \(i\) is the current stage index. The
main routine is \texttt{step(offroad\_rate, crash\_rate, avg\_reward)}, called from
\texttt{main\_train.py} after each episode; it returns a pair
\((\text{target\_npcs},\, \text{event} \in \{\texttt{none},\, \texttt{advance},\, \texttt{rollback}\})\).
The training loop reacts to the event by updating \texttt{CarlaEnv.num\_npc\_vehicles} in place,
so no reconnection to CARLA is required (Section~\ref{sec:design-curriculum}).

\section{Training and evaluation pipelines}\label{sec:impl-pipelines}

Both entry points share the function \texttt{build\_env(args)} that constructs the wrapper chain
of Figure~\ref{fig:wrapper-chain} with the shield type, reward weights, curriculum NPC count and
map chosen by the command line. Keeping the same builder on both sides guarantees that the
evaluation environment matches the training environment exactly.

\subsection{Training pipeline}\label{sec:impl-training}

Figure~\ref{fig:train-algorithm} summarises the training loop.

\begin{figure}[h]
  \centering
  \includegraphics[width=0.95\linewidth]{imgs/code/training_loop}
  \caption{High-level training loop (\texttt{main\_train.py}). The inner loop collects
           transitions and appends the \emph{executed} action (after shield correction) to the
           rollout buffer. A PPO update is triggered every 2048 steps. After each episode, the
           curriculum manager is polled; if it advances or rolls back a stage, the NPC count
           is updated in place without restarting the environment.}
  \label{fig:train-algorithm}
\end{figure}

The key correctness invariant of the training loop is step~11 of Figure~\ref{fig:train-algorithm}:
the log probability stored in the buffer corresponds to the \emph{executed} action
\(a^\star_t\), not the sampled \(a_t\). This is what makes the PPO gradient unbiased in the
presence of shield interventions (Section~\ref{sec:impl-shield-credit}).

\subsection{Evaluation pipeline}\label{sec:impl-evaluation}

The evaluation pipeline (\texttt{main\_eval.py}) rebuilds the wrapper chain with the same
\texttt{build\_env} function but sets all reward weights to zero so that the reported episode
reward is the pure CARLA base reward. The shield is still attached so that action modifications
are visible. An optional Matplotlib dashboard displays the LIDAR polar plot, speed gauge, lateral
offset marker, and a monospaced panel with live information (risk level, shield activation,
steering, throttle/brake, total distance, collisions). At the end of the run a
\texttt{SafetyMetricsReporter} prints a summary table aggregating success rate, crash rate,
off-road rate, shield activations and LIDAR-derived minimum distances.

\section{Persistence and reproducibility}\label{sec:impl-repro}

\subsection{Checkpoint format}\label{sec:impl-checkpoints}

Checkpoints are saved with \texttt{torch.save} as a dictionary of two entries: the
\texttt{state\_dict} of the \texttt{ActorCritic} network and the state of the
\texttt{RunningMeanStd} normaliser. The normaliser state is essential for evaluation: if it were
missing, the running statistics would be rebuilt from scratch and the input distribution would not
match the one the policy trained on. A shape check on load tolerates the legacy format (state
dictionary only) for backwards compatibility with the earliest checkpoints.

\subsection{Seeds and determinism}\label{sec:impl-seeds}

Two distinct seeds are used:

\begin{itemize}
  \item \texttt{seed=42} for training. Propagated to NumPy, Python's \texttt{random} module,
        PyTorch, CARLA's Traffic Manager and the environment's episode-counter random state.
  \item \texttt{seed=100} for evaluation, intentionally different from the training seed so
        that evaluation trajectories are not a subset of those seen during learning.
\end{itemize}

Full bit-exact determinism across hardware is not achievable because of PyTorch's cuDNN
non-reproducibility and CARLA's floating-point physics divergence. However, determinism is
sufficient for the first fifty episodes of a run, which is long enough to diagnose regressions
introduced by a code change (acceptance criterion~AC-09, Table~\ref{tab:acceptance-criteria}).

\subsection{Tooling}\label{sec:impl-tooling}

The source code is linted and formatted with \texttt{ruff}. Reward regressions are caught by
\texttt{pytest tests/test\_reward\_balance.py}. The utility scripts under \texttt{utils/}
provide helpers to list, inspect and compare checkpoints (\texttt{ModelManager}), to analyse
the metrics of finished runs (\texttt{ExperimentAnalyzer}) and to generate configuration files
from a template (\texttt{ConfigurationTemplate}). None of these utilities are required for
training; they are purely auxiliary.
